% SEGMENT OLP-0087-S01
% Part: first-order-logic
% Chapter: natural-deduction
% Section: quantifier-rules

\documentclass[../../../include/open-logic-section]{subfiles}

\begin{document}

\olfileid{fol}{ntd}{qrl}

\olsection{அளவையடை விதிகள்}

\subsection{$\lforall$ க்கான விதிகள்}

\begin{defish}
\AxiomC{$!A(a)$}
\RightLabel{\Intro{\lforall}}
\UnaryInfC{$\lforall[x][\Atom{!A}{x}]$}
\DisplayProof
\hfill
\AxiomC{$\lforall[x][\Atom{!A}{x}]$}
\RightLabel{\Elim{\lforall}}
\UnaryInfC{$!A(t)$}
\DisplayProof
\end{defish}

$\lforall$ க்கான விதிகளில், $t$ ஒரு மூடிய சொல் (எந்த மாறியையும் கொண்டிராத
சொல்); $a$ ஆனது முடிவு~$\lforall[x][!A(x)]$ இலும், முன்கோள்~$!A(a)$ உடன்
முடியும் !!{derivation}[loc] !!{undischarged} ஆக இருக்கும் எந்த எடுகோளிலும்
இடம்பெறாத !!a{constant}. $a$ ஐ \Intro{\lforall} அனுமானத்தின்
\emph{தனிமாறி} என்கிறோம்.\footnote{மேலுள்ள விதியில் $a$ ஒரு மாறிலிக்
குறியாக இருந்தபோதும், ``தனிமாறி'' என்ற சொல்லைப் பயன்படுத்துகிறோம். இதற்கு
வரலாற்றுக் காரணங்கள் உள்ளன.}

% SEGMENT OLP-0087-S02
\subsection{$\lexists$ க்கான விதிகள்}

\begin{defish}
\AxiomC{$\Atom{!A}{t}$}
\RightLabel{\Intro{\lexists}}
\UnaryInfC{$\lexists[x][\Atom{!A}{x}]$}
\DisplayProof
\hfill
\AxiomC{$\lexists[x][\Atom{!A}{x}]$}
\AxiomC{[$\Atom{!A}{a}$]$^n$}
\DeduceC{$!C$}
\DischargeRule{\Elim{\lexists}}{n}
\BinaryInfC{$!C$}
\DisplayProof
\end{defish}

% SEGMENT OLP-0087-S03
இங்கும் $t$ ஒரு மூடிய சொல். மேலும் $a$ ஆனது முன்கோள்
$\lexists[x][!A(x)]$ இலும், முடிவு~$!C$ இலும், இரு முன்கோள்களுடன் முடியும்
!!{derivation}s[loc] !!{undischarged} ஆக இருக்கும் எந்த எடுகோளிலும்
($!A(a)$ என்ற எடுகோள்களைத் தவிர) இடம்பெறாத !!a{constant}. $a$ ஐ
\Elim{\lexists} அனுமானத்தின் \emph{தனிமாறி} என்கிறோம்.

% SEGMENT OLP-0087-S04
\Intro{\lforall} அல்லது \Elim{\lexists} அனுமானத்திற்கான முன்கோள்களுக்கு
இட்டுச் செல்லும் !!{derivation}s[loc], தனிமாறி முன்கோள்களிலோ எந்த
!!{undischarged} எடுகோளிலோ இடம்பெறக்கூடாது என்ற நிபந்தனை
\emph{தனிமாறி நிபந்தனை} எனப்படுகிறது.

% SEGMENT OLP-0087-S05
\begin{explain}
$!A$ என்பது !!{variable}~$x$ ஐ கட்டற்றதாகக் கொண்ட !!a{formula} எனில், இதனை
$!A(x)$ என்று எழுதிக் காட்டும் மரபை நினைவுகூர்க. அதே சூழலில், $!A(t)$ என்பது
$\Subst{!A}{t}{x}$ என்பதன் சுருக்கம். எனவே $\Intro\lexists$ விதியைப்
பின்வருமாறும் எழுதலாம்:
\begin{prooftree}
\AxiomC{$\Subst{!A}{t}{x}$}
\RightLabel{\Intro{\lexists}}
\UnaryInfC{$\lexists[x][!A]$}
\end{prooftree}
$t$ ஏற்கெனவே~$!A$ இல் இடம்பெறலாம் என்பதை கவனிக்கவும்; எடுத்துக்காட்டாக,
$!A$ ஆனது~$\Atom{\Obj P}{t,x}$ ஆக இருக்கலாம். ஆகவே,
$\Atom{\Obj P}{t,t}$ இலிருந்து $\lexists[x][\Atom{\Obj P}{t,x}]$ ஐ
அனுமானிப்பது $\Intro\lexists$ இன் சரியான பயன்பாடு—முன்கோளில் வரும் $t$ இன்
ஒன்று அல்லது அதற்கு மேற்பட்ட இடம்பெறுதல்களை, கட்டுண்ட !!{variable}~$x$ ஆல்
``மாற்றலாம்''; அவை அனைத்தையும் மாற்ற வேண்டியதில்லை. எனினும்,
$\Intro\lforall$, $\Elim\lexists$ ஆகியவற்றின் தனிமாறி நிபந்தனைகள்,
!!{constant}~$a$ ஆனது~$!A$ இல் இடம்பெறக்கூடாது என்று கோருகின்றன. எனவே,
$\Atom{\Obj P}{a,a}$ இலிருந்து $\lforall[x][\Atom{\Obj P}{a,x}]$ ஐ
$\Intro\lforall$ பயன்படுத்தி சரியாக அனுமானிக்க முடியாது.
\end{explain}

% SEGMENT OLP-0087-S06
\begin{explain}
\Intro{\lexists}, \Elim{\lforall} ஆகியவற்றில் கட்டுப்பாடுகள் இல்லை;
சொல்~$t$ எதுவாகவும் இருக்கலாம். ஆகவே எந்த நிபந்தனையையும் பற்றி கவலைப்பட
வேண்டியதில்லை. மறுபுறம், \Elim{\lexists}, \Intro{\lforall} விதிகளில்,
!!{constant}~$a$ முடிவிலோ ஓர் !!{undischarged} எடுகோளிலோ எங்கும்
இடம்பெறக்கூடாது என்று தனிமாறி நிபந்தனை கோருகிறது. முறைமை சரித்தன்மை
உடையதாக, அதாவது எந்த !!{undischarged} எடுகோள்களிலிருந்து !!{sentence}s
பின்தொடர்கின்றனவோ அவற்றிலிருந்து மட்டுமே அவற்றை !!{derive} செய்வதாக இருப்பதை
உறுதிசெய்ய இந்நிபந்தனை தேவை. இந்நிபந்தனை இல்லையெனில் பின்வருவது
அனுமதிக்கப்படும்:
\begin{prooftree}
  \AxiomC{$\lexists[x][!A(x)]$}
  \AxiomC{$\Discharge{!A(a)}{1}$}
  \RightLabel{*\Intro{\lforall}}
  \UnaryInfC{$\lforall[x][!A(x)]$}
  \RightLabel{\Elim{\lexists}}
  \BinaryInfC{$\lforall[x][!A(x)]$}
\end{prooftree}
ஆனால் $\lexists[x][!A(x)] \Entails/ \lforall[x][!A(x)]$.

அளவையடைகளுக்கான நீக்கல் விதிகள், !!{variable}s[dat] பதிலாக மூடிய சொற்களை
மட்டுமே இட அனுமதிப்பதால், !!{sentence}s[gen] ஒரு கணத்திலிருந்து வருவிக்கப்படும்
எந்த !!{formula} உம் தானே !!a{sentence} என்பது பின்தொடர்கிறது.
\end{explain}


\end{document}
